\documentclass[conference]{IEEEtran}

\usepackage{amsmath,amssymb}
\usepackage{bm}
\usepackage{graphicx}
\usepackage{cite}
\usepackage{hyperref}
\usepackage{algorithm}
\usepackage{algorithmic}
\usepackage{booktabs}
\usepackage{float}

\title{Interference-Based Quantum Classification of Histopathology Images Using Linear State Overlap}

\author{
\IEEEauthorblockN{
Tarakeshwaran S\IEEEauthorrefmark{1},
Srivaishnav S\IEEEauthorrefmark{1},
Pranesh S\IEEEauthorrefmark{1},
Dr. Gunasundari S\IEEEauthorrefmark{1},
}
\IEEEauthorblockA{\IEEEauthorrefmark{1}
Department of Computer Science and Engineering\\
Velammal Engineering College, Chennai, India\\
tarakeshwaran.sampath@gmail.com\\
vaishnav90100@gmail.com\\
praneshjd2004@gmail.com\\
gunasundari@velammal.edu.in
}
}

\begin{document}
\maketitle

\begin{abstract}
Near-term quantum machine learning (QML) is constrained by two practical bottlenecks: encoding high-dimensional medical images into quantum states and performing measurement-heavy inference when decisions rely on quadratic fidelities. We address this gap with a hybrid medical-QML pipeline in which a classical CNN backbone with a compact neck compresses 96\,$\times$\,96 histopathology patches into a low-dimensional embedding that is amplitude-encoded, followed by a fixed interference circuit that classifies using the \emph{sign of the real linear overlap}. Learning is fully classical: class-representative prototype states are constructed by centroid aggregation of embeddings and compiled into fixed state-preparation unitaries, enabling inference with one fixed circuit per sample and a small constant number of shots. Experimental evaluation uses the PatchCamelyon dataset with statevector simulation and realistic noise injection. On the PCam validation subset, the proposed IQC achieves 88.07\% accuracy while maintaining measurement-efficient, low-shot inference. These results indicate a practical NISQ-compatible classifier tailored for medical image inference under resource constraints.
\end{abstract}

\begin{IEEEkeywords}
NISQ Devices, Quantum Machine Learning, Quantum Classification, Quantum Interference, Linear Overlap, Hybrid Quantum--Classical Systems, Medical Imaging, Histopathology
\end{IEEEkeywords}

\section{Introduction}

\subsection{Background of Quantum Machine Learning}
Quantum machine learning (QML) promises new modeling primitives, but near-term classifiers remain constrained by either unstable variational training or measurement-heavy similarity estimation \cite{benedetti2019parameterized,mcclean2018barren,schuld2019feature,havlicek2019supervised}. These challenges intensify in medical imaging, where raw histopathology patches are high-dimensional and direct quantum encoding of 96\,$\times$\,96 images is infeasible on NISQ hardware.

\subsection{Limitations of Existing QML Approaches}
Current QML pipelines often rely on parameterized circuits that are difficult to train at scale or on similarity estimation procedures that demand large numbers of measurement shots. In clinical imaging tasks, both training stability and measurement efficiency are critical to feasible deployment.

\subsection{Motivation for Interference-Based Classification}
We address this gap with a hybrid medical-QML architecture. A classical CNN backbone with a compact neck compresses each image into a low-dimensional embedding. The embedding is then amplitude-encoded into a quantum state and processed by a fixed interference circuit. Classification is performed by the sign of the real linear overlap between the test state and a class prototype, requiring a small constant number of shots per inference using one fixed circuit per sample. Learning is entirely classical: class-representative prototype states are constructed by centroid aggregation of embeddings, and the quantum processor is used only for inference.

\subsection{Contributions of This Work}
Our contributions are:
\begin{itemize}
    \item A \textbf{hybrid medical-QML pipeline} for lymph node histopathology that makes quantum encoding feasible by using a CNN neck for dimensionality reduction.
    \item A \textbf{linear-interference decision primitive} based on $\mathrm{sign}(\mathrm{Re}\langle\phi|\psi\rangle)$, yielding a hyperplane classifier in Hilbert space.
    \item A \textbf{fixed interference circuit} requiring measurement-efficient inference with a small constant number of shots and no variational training.
    \item A \textbf{fully classical learning phase} using centroid prototypes, decoupling training from quantum execution.
\end{itemize}

\subsection{Structure of the Paper}
The remainder of this paper is organized as follows. Section~II reviews the relevant QML literature. Section~III summarizes existing system paradigms and their limitations. Section~IV presents the proposed interference-based classifier, including the hybrid pipeline, encoding, and decision mechanism. Section~V details the implementation setup, datasets, and evaluation protocol. Section~VI discusses measurement efficiency, noise robustness, limitations, and scalability. Section~VII concludes and outlines future work.

\section{Literature Review}

Quantum machine learning (QML) has progressed from early linear-algebraic subroutines to hardware-focused models tailored to NISQ constraints. We summarize the main approaches and position the linear-interference framework within the remaining gaps.

\subsection{Variational Quantum Classifiers}

Variational quantum classifiers are the most well-known type of machine learning model from the NISQ era \cite{benedetti2019parameterized,abbas2021power}. VQCs use a parameterized quantum circuit (PQC) as an ansatz, and they change the parameters to make a classical cost function as small as possible. VQCs are very adaptable, but they are also very hard to train. Gradient-based optimization is often not possible for larger systems because there are empty plateaus in the parameter space where gradients disappear exponentially with the number of qubits \cite{mcclean2018barren}. Also, VQCs need a lot of measurement shots at each optimization step to get low-variance gradient estimates, which puts a lot of strain on the quantum processor's duty cycle.

\subsection{Quantum Kernel and Similarity-Based Methods}

Quantum kernel methods circumvent the optimization challenges associated with VQCs by projecting data into high-dimensional Hilbert space and calculating the similarities between pairs of data points \cite{schuld2019feature,havlicek2019supervised}. These similarities, commonly characterized as the fidelity $F(\rho, \sigma) = \text{Tr}(\sqrt{\sqrt{\rho}\sigma\sqrt{\rho}})^2$, serve as kernel entries in classical algorithms such as Support Vector Machines (SVMs). Kernel approaches are attractive but can be measurement-intensive on quantum hardware. To classify a new point against a dataset of size $N$, one typically requires $O(N)$ similarity measurements. Fidelity is also a quadratic observable that discards relative phase information, which can matter when the decision boundary depends on sign structure in the amplitudes.

\subsection{Quantum State Discrimination and Metric Learning}

Quantum state discrimination (QSD) methodologies regard classification as the process of determining which of several known quantum states a particular system occupies \cite{sergioli2019binary,gonzalez2021classification}. From this point of view, the focus moves from circuit parameters to the way states are arranged in Hilbert space. Quantum metric learning, on the other hand, tries to find a feature map that makes the distance between states that represent different classes as large as possible and the distance between states that represent the same class as small as possible. Our research is based on these geometric principles, but we add a new way of making inferences: we use a fixed interference circuit to get a linear overlap signal, which is a simpler and more efficient way to measure than full state tomography or fidelity estimation.

\subsection{Interference-Based Quantum Classification}

Works like the distance-based classifier by Schuld et al. \cite{schuld2017distance,schuld2018book} pioneered the use of the Hadamard test for pattern recognition by calculating distances between vectors, and kernel-based variants compute signed fidelity sums in superposition \cite{blank2020kernel,blank2022compact}. Our work differs operationally in two ways: (i) the decision variable is the \textbf{sign of the real overlap}, not a fidelity or distance estimate, and (ii) learning is a fully classical centroid construction, with a \textbf{fixed interference circuit} used only at inference. This yields a measurement-efficient primitive that avoids variational training and reduces circuit executions per sample.

\begin{table*}[t]
\centering
\caption{Comparison of Quantum Classification Paradigms}
\label{tab:comparison}
\begin{tabular}{@{}lllll@{}}
\toprule
Feature & Variational (VQC) & Kernel (QSVM) & Fidelity-Based & \textbf{Linear Interference (IQC)} \\ \midrule
Decision Basis & Probabilities $P(y|x)$ & High-dim Kernels & Quadratic Overlap $F$ & \textbf{Linear Overlap $\mathrm{Re}\langle \phi|\psi \rangle$} \\
Learning Mode & In-circuit Optimization & Classical SVM ($O(N^2)$) & Static/Adaptive & \textbf{Static Classical Aggregation} \\
Quantum Circuit & Parameterized (Trainable) & Fixed (Feature Map) & Fixed (Swap Test) & \textbf{Fixed (Hadamard Test)} \\
Shot Complexity & High ($O(M \cdot N)$) & Moderate ($O(N)$) & Moderate ($O(1/\epsilon^2)$) & \textbf{Low (constant shots)} \\
Phase Sensitivity & No & No & No & \textbf{Retains sign of real interference} \\
Optimization Risk & Barren Plateaus & None (in quantum) & None & \textbf{None} \\ \bottomrule
\end{tabular}
\end{table*}

\subsection{Challenges in Existing Models}
Existing models exhibit three recurring limitations: (i) variational training suffers from barren plateaus and high-variance gradient estimates, (ii) kernel and similarity-based approaches incur heavy measurement cost due to repeated overlap estimation, and (iii) fidelity-based similarity discards the sign of interference, which can be informative for decision boundaries. These issues motivate alternatives that preserve interference structure while reducing measurement overhead.

\section{Existing System}

\subsection{Variational Quantum Classifiers}
VQCs use parameterized quantum circuits optimized to minimize a classical loss. They provide flexible hypothesis classes but face optimization challenges such as barren plateaus and high shot cost for gradient estimation. The dependence on iterative quantum-classical feedback makes them hardware-intensive.

\subsection{Quantum Kernel-Based Classifiers}
Quantum kernel methods map data into high-dimensional Hilbert spaces and compute pairwise similarities for classical SVMs. While they avoid in-circuit training, inference typically scales with dataset size due to the need to estimate kernel values, which can be measurement-intensive.

\subsection{Fidelity-Based Quantum Similarity Methods}
Fidelity-based schemes assess similarity using $F(\phi, \psi) = |\langle \phi | \psi \rangle|^2$. This quadratic observable is robust but discards relative phase, leading to curved decision boundaries and higher shot requirements for precise estimation.

\subsection{Limitations of Existing Systems}
Across these paradigms, the main inefficiencies arise from optimization instability (VQCs), large measurement budgets (kernel and fidelity methods), and loss of sign information in quadratic overlaps. These limitations motivate a linear-interference decision rule that retains phase-sensitive information while using a fixed, low-shot circuit.

\section{Proposed Work}

\subsection{Overview of Hybrid Pipeline}
The proposed Interference-Based Quantum Classifier (IQC) follows a modular pipeline: \emph{Input image} $\rightarrow$ \emph{CNN backbone} $\rightarrow$ \emph{compact neck embedding} $\rightarrow$ \emph{quantum state preparation} $\rightarrow$ \emph{interference-based decision}. The classical module compresses high-dimensional images into an embedding compatible with a small qubit register, while the quantum module performs a fixed interference measurement to obtain a sign-based decision.

Training is fully classical. We first train the CNN backbone and neck on PCam using standard cross-entropy loss to produce discriminative embeddings. After training, the CNN+neck is frozen and used to generate embeddings for all training samples. These embeddings are aggregated into class prototypes (Section~\ref{sec:centroids}), which are compiled into fixed state-preparation unitaries. Quantum hardware is used only for inference: for each test sample we compute its embedding, amplitude-encode it as $|\psi\rangle$, run the fixed interference circuit against the class prototype, and assign a label based on the sign of the interference score estimated over a small constant number of shots.

Fig.~\ref{fig:flowchart} illustrates the end-to-end pipeline.
\begin{figure}[ht]
    \centering
    \includegraphics[width=1.02\linewidth]{flowchart_fig1.png}
    \caption{Flowchart of the hybrid classical-quantum inference pipeline using the Interference-Based Quantum Classifier (IQC).}
    \label{fig:flowchart}
\end{figure}

\subsection{Classical Feature Extraction Module}
The hybrid pipeline begins with a classical CNN backbone that processes 96\,$\times$\,96 RGB histopathology patches and extracts hierarchical features correlated with malignancy. To make quantum encoding feasible, we insert a compact ``neck'' module that compresses the feature maps into a low-dimensional embedding. The neck acts as a learned projection (analogous to a task-specific PCA), implemented with global average pooling followed by a fully connected layer. The embedding dimension $d$ is chosen to match the available qubit budget via $n = \lceil \log_2 d \rceil$. A visual representation of the backbone and feature neck is shown in Fig.~\ref{fig:architecture}.

\begin{figure*}[t]
\centering
\includegraphics[width=1.0\linewidth]{{architecture diagram}.png}
\caption{Hybrid pipeline: image $\rightarrow$ CNN backbone $\rightarrow$ neck (dimensionality reduction) $\rightarrow$ quantum encoding $\rightarrow$ interference-based classifier tail. The quantum stage uses $n$ data qubits and one ancilla for sign-based inference (small constant shots).}
\label{fig:architecture}
\end{figure*}

\subsubsection{Module-wise Tensor Shapes}
The end-to-end data flow is:
\begin{itemize}
    \item Input image: $(96, 96, 3)$
    \item CNN feature maps: $(h, w, c)$ (architecture-dependent)
    \item Global average pooling: $(c)$
    \item Neck projection: $(d)$
    \item Quantum feature vector (normalized and zero-padded if needed): $(2^n)$
    \item Quantum register: $n$ data qubits + 1 ancilla
    \item Output: 1 interference score (sign used for label)
\end{itemize}

This structure explicitly defines the image $\rightarrow$ neck $\rightarrow$ quantum classifier tail, ensuring the quantum stage operates on a manageable number of qubits while preserving discriminative features.

\subsection{Quantum Encoding Module}
A critical component of the IQC is the explicit mapping of classical feature embeddings into quantum states. We adopt \textbf{amplitude encoding} as the primary feature map $\Phi: \mathbb{R}^d \rightarrow \mathbb{C}^{2^n}$. A classical embedding $\mathbf{x} \in \mathbb{R}^d$ is first normalized and (if needed) zero-padded to length $2^n$, then mapped to an $n$-qubit state via a state-preparation unitary $U(\mathbf{x})$:

\begin{equation}
|\psi\rangle = \sum_{i=0}^{2^n - 1} x_i |i\rangle,
\end{equation}

with $n = \lceil \log_2 d \rceil$. Amplitude encoding provides the highest information density, encoding $d$ features into $n$ qubits. In our implementation, amplitude encoding is realized via Qiskit's \texttt{StatePreparation} gate, which constructs a sequence of uniformly controlled rotations to prepare any normalized complex vector $\mathbf{x}$ as a quantum state $|\psi\rangle$. The interference mechanism itself is independent of the specific compilation used.

We consider three primary encoding strategies adapted for the IQC framework:
\begin{enumerate}
    \item \textbf{Amplitude Encoding}: Maps a normalized vector $\mathbf{x} \in \mathbb{R}^d$ to the amplitudes of a $2^n$-dimensional quantum state. This provides the highest information density ($O(2^n)$ features) but requires complex state-preparation circuits. In IQC, amplitude encoding ensures that the linear overlap in Hilbert space is directly proportional to the Euclidean dot product in the input space.
    \item \textbf{Angle Encoding}: Maps each feature $x_j$ to the rotation angle of a qubit, $|\psi\rangle = \bigotimes_j ( \cos x_j |0\rangle + \sin x_j |1\rangle)$. This is hardware-efficient but introduces nonlinearities (trigonometric functions) into the decision signal.
    \item \textbf{Basis Encoding}: Maps classical bitstrings directly to computational basis states. While simple, it does not exploit the superposition-based interference that defines the IQC's advantage.
\end{enumerate}
In this work, we use amplitude encoding for all reported experiments.

\subsubsection{Linear Separability in High-Dimensional Feature Spaces}
According to Cover's Theorem \cite{cover1965geometrical}, a complex pattern classification problem cast in a high-dimensional space nonlinearly is more likely to be linearly separable than in a low-dimensional space. By mapping data into the $2^n$-dimensional Hilbert space, the IQC leverages this high dimensionality to separate classes using a simple hyperplane.

However, unlike kernel methods that use an implicit feature map, the IQC works with an \textbf{explicit feature map} $\Phi(\mathbf{x})$. The linear overlap $\mathrm{Re}\langle \phi | \psi \rangle$ is computed directly in the state space. The success of the static centroid construction (Section~\ref{sec:centroids}) depends on whether the chosen $\Phi$ places the class distributions in convex regions that can be separated by a single hyperplane.

\subsection{Linear Interference Decision Mechanism}
The central theoretical contribution of this work is the application of linear interference as the main signal for quantum classification.

\subsubsection{Linear vs. Quadratic Similarity}
In conventional schemes, similarity is often expressed through fidelity $F(\phi, \psi) = |\langle \phi | \psi \rangle|^2$. Although fidelity is a natural measure of state similarity, its quadratic form leads to three issues in NISQ inference:
\begin{enumerate}
    \item \textbf{Sign Loss}: The absolute square operation $|\cdot|^2$ discards the sign of the real interference term.
    \item \textbf{Non-linear Decision Boundaries}: Fidelity-induced decision boundaries are second-order surfaces in amplitude space.
    \item \textbf{Shot Noise Sensitivity}: Estimating fidelity to high precision requires $O(1/\epsilon^2)$ measurement shots.
\end{enumerate}

In contrast, the linear overlap $\langle \phi | \psi \rangle$ is complex-valued. We focus on its real part:
\begin{equation}
s(\psi) = \mathrm{Re}\langle \phi | \psi \rangle.
\end{equation}

\subsubsection{Decision Function and Labeling Rule}
We define the binary classification decision based on the sign of the real overlap between the input state $|\psi\rangle$ and a predetermined class-representative state $|\phi\rangle$. The predicted label $\hat{y}$ is given by:
\begin{equation}
\hat{y} = \text{sign}\left( \mathrm{Re}\langle \phi | \psi \rangle \right).
\end{equation}
Setting the threshold at zero corresponds to a balanced classification. If prior class probabilities are known, a bias term $b \in \mathbb{R}$ can be added classically:
\begin{equation}
\hat{y} = \begin{cases} +1 & \mathrm{Re}\langle \phi | \psi \rangle + b \geq 0 \\ -1 & \mathrm{Re}\langle \phi | \psi \rangle + b < 0 \end{cases}
\end{equation}

\subsubsection{Decision Geometry: The Hilbert Space Hyperplane}
Geometrically, the decision rule $\mathrm{Re}\langle \phi | \psi \rangle = 0$ defines a \textbf{linear hyperplane} that partitions the $2^n$-dimensional Hilbert space into two half-spaces. Fig.~\ref{fig:geometry_comparison} contrasts this linear-interference boundary with the curved, quadratic boundary induced by fidelity-based similarity, and Fig.~\ref{fig:fidelity_geometry} shows the fidelity geometry.

To see this, let $|\psi\rangle$ and $|\phi\rangle$ be represented by their amplitude vectors $\mathbf{c}_\psi, \mathbf{c}_\phi \in \mathbb{C}^{2^n}$. The real part of the inner product can be written as:
\begin{equation}
\mathrm{Re}\langle \phi | \psi \rangle = \frac{1}{2}\left( \langle \phi | \psi \rangle + \langle \psi | \phi \rangle \right) = \text{Re}(\mathbf{c}_\phi^\dagger \mathbf{c}_\psi).
\end{equation}

\begin{figure}[t]
\centering
\includegraphics[width=1.0\linewidth]{{Interference-Based decision Geometery}.png}
\caption{Comparison of decision boundaries: Linear Interference (left) vs. Fidelity-based Quadratic (right). The interference rule $\mathrm{Re}\langle \phi | \psi \rangle = 0$ induces a linear hyperplane and retains the sign of the real interference term, while fidelity-based rules result in curved boundaries that discard sign information.}
\label{fig:geometry_comparison}
\end{figure}

\begin{figure}[H]
\centering
\includegraphics[width=0.8\linewidth]{{Fidelity based similarity geometry}.png}
\caption{Detailed geometric profile of fidelity-based similarity measures in Hilbert space, illustrating the nonlinear curvature of the decision surface.}
\label{fig:fidelity_geometry}
\end{figure}

\subsubsection{Formal Derivation of the Decision Score}
Consider the normalized state $|\psi\rangle = \sum_i a_i |i\rangle$ and the representative state $|\phi\rangle = \sum_i b_i |i\rangle$. The interference term $\text{Re}\langle \phi | \psi \rangle$ can be expanded in terms of the real and imaginary parts of the coefficients:
\begin{equation}
    \langle \phi | \psi \rangle = \sum_i b_i^* a_i = \sum_i (\text{Re } b_i - i \text{Im } b_i)(\text{Re } a_i + i \text{Im } a_i)
\end{equation}
Taking the real part, we obtain:
\begin{equation}
    \text{Re}\langle \phi | \psi \rangle = \sum_i (\text{Re } a_i \text{Re } b_i + \text{Im } a_i \text{Im } b_i)
\end{equation}
This summation is identical to the Euclidean inner product in $\mathbb{R}^{2 \cdot 2^n}$. The decision boundary $\text{Re}\langle \phi | \psi \rangle = 0$ is therefore the locus of all states in $\mathcal{H}$ that are "mutually orthogonal" under this real-part inner product.

\subsection{Quantum Interference Circuit}
The accessibility of the linear overlap signal $\mathrm{Re}\langle \phi | \psi \rangle$ is the defining hardware requirement of our framework. Since standard projective measurements $\langle \psi | O | \psi \rangle$ are quadratic in the state amplitudes, we utilize an ancilla-assisted interference circuit. In this framework, one inference corresponds to executing a fixed interference circuit on an input state and assigning a label based on the sign of the estimated overlap.

\subsubsection{Physical Realization via the Hadamard Test}
As with all gate-based quantum classifiers, state preparation is treated as a separate cost; the inference complexity discussed here refers exclusively to the decision observable and measurement process.

The real part of the overlap between two quantum states can be directly extracted using the Hadamard test \cite{ekert2002direct}. The circuit requires one additional ancilla qubit and the original register used to store the data states.

The procedure for evaluating $\mathrm{Re}\langle \phi | \psi \rangle$ is as follows:
\begin{enumerate}
    \item \textbf{Initialization}: Prepare the ancilla qubit in $|0\rangle$ and the register in a state $|\psi\rangle$.
    \item \textbf{First Hadamard}: Apply a Hadamard gate $H$ to the ancilla qubit, putting it in the state $\frac{1}{\sqrt{2}}(|0\rangle + |1\rangle)$.
    \item \textbf{Controlled Operation}: Apply a controlled unitary $U$ that transforms $|\psi\rangle$ toward $|\phi\rangle$. In the most hardware-efficient setup, if we can prepare both $|\psi\rangle$ and $|\phi\rangle$ via unitaries $V_\psi$ and $V_\phi$ from $|0\rangle$, the operation becomes a controlled version of $V_\phi V_\psi^\dagger$.
    \item \textbf{Second Hadamard}: Apply a second Hadamard gate to the ancilla.
    \item \textbf{Measurement}: Measure the ancilla qubit in the computational basis.
\end{enumerate}

The probability of measuring the ancilla in state $|0\rangle$ is:
\begin{equation}
P(0) = \frac{1}{2} + \frac{1}{2} \mathrm{Re}\langle \phi | \psi \rangle.
\end{equation}
The interference signal (our decision score) is then obtained from the expectation value of the ancilla's $Z$-observable:
\begin{equation}
\langle Z_{\text{anc}} \rangle = P(0) - P(1) = \mathrm{Re}\langle \phi | \psi \rangle.
\end{equation}

\begin{figure}[t]
\centering
\includegraphics[width=1.0\linewidth]{transition_isdo.png}
\caption{Circuit diagram of the ISDO-B implementation. The controlled transition unitary $U_{\chi\psi}$ is used to estimate the linear overlap $\mathrm{Re}\langle \phi | \psi \rangle$ by measuring the ancilla qubit expectation value over a small constant number of shots.}
\label{fig:circuit}
\end{figure}

\subsubsection{ISDO-B: Hardware-Efficient Linear Interference}
While the canonical Hadamard test relies on controlled state preparation for both $|\psi\rangle$ and $|\phi\rangle$, the \textbf{ISDO-B (Interference-Based Static Decision Observable)} variant utilizes a controlled transition unitary $U_{\chi\psi}$ where $U_{\chi\psi}|\psi\rangle = |\phi\rangle$. We assume $U_{\chi\psi}$ is a state-transfer operator acting such that the target subspace is correctly mapped; its action on the orthogonal subspace is unconstrained as it does not contribute to the interference signal.

The circuit follows a modified Hadamard test structure:
\begin{enumerate}
    \item Initialize ancilla in $|0\rangle$ and data register in $|0\rangle^{\otimes n}$.
    \item Prepare test state $|\psi\rangle$ on the data register.
    \item Apply Hadamard gate to the ancilla: $\frac{1}{\sqrt{2}}(|0\rangle+|1\rangle) \otimes |\psi\rangle$.
    \item Apply controlled transition unitary $CU_{\chi\psi}$ targeting the data register: $\frac{1}{\sqrt{2}}(|0\rangle |\psi\rangle + |1\rangle |\phi\rangle)$.
    \item Apply Hadamard gate to the ancilla and measure in the $Z$-basis.
\end{enumerate}
The expectation value of the ancilla measurement yields $\langle Z_{\text{anc}} \rangle = \mathrm{Re}\langle \phi | \psi \rangle$. This formulation is particularly advantageous in static regimes where the class prototype $|\phi\rangle$ is precomputed, allowing for the engineering of specialized observables that avoid the overhead of full controlled-state preparation of arbitrary centroids.

\subsubsection{Single-Circuit Inference Cost}
The IQC uses a single fixed interference circuit per sample and reads out one ancilla qubit to obtain the decision sign. In our experiments, a small constant number of shots is sufficient to stabilize the sign decision, in contrast to fidelity-based methods that require high-precision probability estimation.

\subsection{Class Representative State Construction}
\label{sec:centroids}
A pivotal design choice in the IQC framework is the complete decoupling of the learning process from quantum execution. Unlike VQCs, where learning involves an iterative quantum-classical loop, our framework performs all structural learning classically before anything is loaded onto the quantum processor.

\subsubsection{Centroid-Based State Construction}
The learning objective is to find a normalized quantum state $|\phi^{(c)}\rangle$ that best represents the distribution of training samples belonging to class $c \in \{-1, +1\}$. In this work, we adopt a \textbf{centroid-based aggregation rule}, which is conceptually equivalent to calculating the mean vector in Hilbert space. This centroid-based approach is functionally analogous to classical Linear Discriminant Analysis (LDA), where class means are used to establish a linear decision boundary, albeit implemented directly within the high-dimensional quantum Hilbert space.

Given a subset of training states $\mathcal{S}_c = \{|\psi_i\rangle : y_i = c\}$, where each $|\psi_i\rangle$ is an amplitude-encoded vector in $\mathbb{C}^d$, the unnormalized class-representative vector is:
\begin{equation}
| \tilde{\phi}^{(c)} \rangle = \sum_{|\psi_i\rangle \in \mathcal{S}_c} w_i |\psi_i\rangle,
\end{equation}
where $w_i$ are importance weights. In the simplest case, where all training samples are considered equally, $w_i = 1/|\mathcal{S}_c|$. The finalized class-representative state is obtained by normalization:
\begin{equation}
|\phi^{(c)}\rangle = \frac{| \tilde{\phi}^{(c)} \rangle}{\| | \tilde{\phi}^{(c)} \rangle \|}.
\end{equation}

Algorithm~\ref{alg:learning} details the classical preprocessing steps required to construct these states.

\begin{algorithm}[ht]
\caption{Static Learning of Class-Representative States}
\label{alg:learning}
\begin{algorithmic}[1]
\REQUIRE Training dataset $\mathcal{D} = \{(\mathbf{x}_i, y_i)\}_{i=1}^N$
\STATE Select feature map $\Phi: \mathbb{R}^d \rightarrow \mathbb{C}^{2^n}$ (e.g., Amplitude Encoding)
\STATE Initialize class accumulators: $\mathbf{v}_+ \leftarrow \mathbf{0}, \mathbf{v}_- \leftarrow \mathbf{0}$
\FOR{each $(\mathbf{x}_i, y_i) \in \mathcal{D}$}
    \STATE Compute state vector: $|\psi_i\rangle \leftarrow \Phi(\mathbf{x}_i)$
    \IF{$y_i = +1$}
        \STATE $\mathbf{v}_+ \leftarrow \mathbf{v}_+ + |\psi_i\rangle$
    \ELSE
        \STATE $\mathbf{v}_- \leftarrow \mathbf{v}_- + |\psi_i\rangle$
    \ENDIF
\ENDFOR
\STATE Normalize: $|\phi^{(+)}\rangle \leftarrow \mathbf{v}_+/\|\mathbf{v}_+\|, |\phi^{(-)}\rangle \leftarrow \mathbf{v}_-/\|\mathbf{v}_-\|$
\RETURN $|\phi^{(+)}\rangle, |\phi^{(-)}\rangle$
\end{algorithmic}
\end{algorithm}

This aggregation is performed entirely in classical memory using the complex amplitude vectors of the data embeddings. For a dataset of size $N$ and a feature space of dimension $d$, this classical step scales as $O(Nd)$, which is significantly faster than the $O(M \cdot N)$ scaling of VQC training (where $M$ is the number of optimization iterations).

\subsubsection{State Selection and Prototype Extraction}
In some applications, a simple average may not be the most representative summary of a class, especially if the class distribution is multimodal. However, within the scope of this work, we assume a single-mode distribution for each class. The state $|\phi^{(c)}\rangle$ essentially acts as a "prototype" or "centroid" in the quantum feature space. The linear interference score $\text{Re}\langle \phi | \psi \rangle$ can then be interpreted as the projection of the test sample $|\psi\rangle$ onto this class prototype.

\subsubsection{Static Nature and Hybrid Efficiency}
The resulting states $|\phi^{(+)}\rangle$ and $|\phi^{(-)}\rangle$ are static. Once computed, they are used to derive the control sequences for the interference circuit described earlier. There are no weight updates, no memory growth, and no backpropagation through quantum gates.

By treating the quantum device as a static inference engine rather than a trainable model, we eliminate several pathologies common in hybrid QML:
\begin{itemize}
    \item \textbf{Gradient Noise}: Classical aggregation is deterministic and does not suffer from the sampling variance of quantum-derived gradients.
    \item \textbf{Barren Plateaus}: We avoid the optimization landscape altogether. The "training" is a simple linear aggregation with a guaranteed closed-form solution.
    \item \textbf{I/O Overhead}: The classical-quantum bridge is traversed only for data loading and the final inference shot, avoiding the high-latency loops of variational training.
\end{itemize}

\section{Implementation}

\subsection{Dataset}
The PatchCamelyon (PCam) dataset is a widely recognized benchmark for histopathological cancer detection, specifically designed for lymph node metastasis identification \cite{pcam_veeling2018,pcam_01,pcam_03}. It consists of 96\,$\times$\,96 RGB image patches extracted from whole-slide lymph node histology scans. Labels indicate whether tumor tissue appears in the central 32\,$\times$\,32 region of the patch, yielding a binary task: metastatic (tumor) vs benign (non-tumor) \cite{pcam_veeling2018}. The full dataset contains 327{,}680 patches and is nearly class-balanced; the standard Kaggle split (duplicates removed) contains approximately 220k training patches and 57.5k test patches \cite{pcam_veeling2018,pcam_03}.

Clinically, accurate detection of metastatic cancer in lymph nodes is crucial for cancer staging and treatment planning. The PCam dataset captures the morphological complexity of lymph node tissue, where malignant regions exhibit distinct cellular structures and tissue organization compared to benign areas. Fig.~\ref{fig:benign_malignant} shows representative examples of benign and malignant patches from the dataset, highlighting key histological differences relevant to classification.

\begin{figure}[ht]
    \centering
    \includegraphics[width=0.45\linewidth]{benign.png}
    \includegraphics[width=0.45\linewidth]{malignant.png}
    \caption{Representative lymph node histopathology patches from the PCam dataset. Left: Benign tissue showing normal lymphatic architecture. Right: Malignant tissue exhibiting disrupted cellular morphology and tumor infiltration.}
    \label{fig:benign_malignant}
\end{figure}

\subsection{CNN Feature Extraction}
The CNN backbone and neck are trained classically using cross-entropy loss on PCam. We follow the original PCam train/val/test split provided by the dataset: the backbone is trained on the \texttt{train} split and validated on the \texttt{val} split for model selection. Due to simulation constraints and the compute-intensive nature of statevector inference, the \texttt{val} split (held out from training) is also used for the final downstream evaluation of the quantum stage, with the \texttt{test} split reserved for future hardware validation.

Training uses a lightweight PCamCNN with three convolutional blocks (3$\times$3 conv + batch normalization + ReLU, with max-pooling after the first two blocks), followed by adaptive average pooling and a linear embedding head of dimension $d=32$ with ReLU activation. A temporary linear classifier is attached during training and discarded after convergence. We use Adam (learning rate $10^{-3}$, weight decay $10^{-4}$), batch size 64, and a ReduceLROnPlateau scheduler (factor 0.5, patience 2) with early stopping (patience 10). Data augmentation is limited to label-preserving transforms: random horizontal/vertical flips and light color jitter, followed by normalization to mean $[0.5,0.5,0.5]$ and std $[0.5,0.5,0.5]$. Validation and embedding extraction use deterministic normalization only.

After training, the CNN+neck (PCamCNN) is frozen and used to produce embeddings for all training samples. These embeddings are then aggregated into class prototypes (Section~\ref{sec:centroids}). The quantum processor is used only at inference.

\subsection{Quantum Encoding Setup}
In our experiments, $d = 32$ and $n = 5$, so the quantum register consists of 5 data qubits and 1 ancilla qubit (6 total). Each embedding $\mathbf{x} \in \mathbb{R}^{32}$ is amplitude-encoded using Qiskit's \texttt{StatePreparation} unitary, which realizes the mapping $|\psi\rangle = \sum_{i=0}^{31} x_i|i\rangle$.

\subsubsection{Depth Analysis of Amplitude Encoding}
While the IQC inference circuit is fixed and short-depth (constant relative to state preparation), the overall performance is bounded by the cost of preparing the states $|\psi\rangle$ and $|\phi\rangle$. For a $d$-dimensional input vector, amplitude encoding into $n = \log_2 d$ qubits requires a unitary $U$ such that $U|0\rangle^{\otimes n} = \sum_{i=1}^d x_i |i\rangle$.

The standard decomposition of such a unitary using the Plesch-Brukner algorithm \cite{plesch2011quantum} results in a circuit with $O(d)$ CNOT gates and $O(d)$ single-qubit rotations. While this is exponential in the number of qubits $n$, it is linear in the number of features $d$. For NISQ deployment, one often uses approximate state preparation techniques or hardware-efficient ans\"atze that trade off exact representation for reduced depth. In our simulation, we consider an exact decomposition but point out that the linear interference mechanism of the IQC is still valid independently of the state preparation technique used, as long as the relative phase information is maintained.

\subsection{Simulation Environment}
Simulations are carried out via a statevector quantum simulator, with a tailored noise model added to simulate NISQ hardware characteristics. The shot complexity is measured relative to the physical measurement step, not simulation time. All results below were obtained in a single experimental run; multi-seed statistical validation is a planned future step.
\begin{itemize}
    \item \textbf{Dataset}: PatchCamelyon (PCam), standard Kaggle split (duplicates removed) \cite{pcam_veeling2018}. Due to statevector simulation memory constraints, we use a balanced held-out subset of \textbf{5000 samples} (2500 per class) drawn from the PCam validation partition.
    \item \textbf{Feature Extraction}: A trained PCamCNN backbone with a compact neck reduces each 96\,$\times$\,96 RGB patch to an embedding of dimension $d = 32$. Embeddings are L2-normalized before quantum encoding.
    \item \textbf{Quantum Encoding}: Each embedding $\mathbf{x} \in \mathbb{R}^{32}$ is amplitude-encoded into $n = \lceil \log_2 32 \rceil = 5$ data qubits plus 1 ancilla (6 qubits total) using Qiskit's \texttt{StatePreparation} unitary, which realizes the mapping $|\psi\rangle = \sum_{i=0}^{31} x_i|i\rangle$.
    \item \textbf{Noise Model}: A depolarizing noise channel with error probability $p \in [0.01, 0.1]$ and thermal relaxation times ($T_1$, $T_2$) consistent with current superconducting processors.
\end{itemize}

\subsection{Benchmark Models}
Our IQC (ISDO-B variant) is compared to Fidelity-SWAP, Linear SVM, Logistic Regression, and $k$-NN ($k=3$) on the same 5000-sample subset.

\subsection{Evaluation Metrics}
Fig.~\ref{fig:accuracy} shows the classification accuracy comparison between the proposed IQC, VQC, and fidelity-based methods across different noise levels.

\begin{figure}[t]
\centering
\includegraphics[width=0.85\linewidth]{accuracy_comparison.png}
\caption{Classification accuracy comparison between the proposed IQC, VQC, and Fidelity-based methods across different noise levels. The IQC maintains stable performance even as noise increases.}
\label{fig:accuracy}
\end{figure}

Table~\ref{tab:metrics} presents classification accuracy at a fixed noise level ($p=0.01$).

\begin{table}[ht]
\centering
\caption{Classification Performance Comparison: Validation of IQC accuracy against quantum and classical baselines under symmetric depolarizing noise ($p=0.01$).}
\label{tab:metrics}
\begin{tabular}{l c}
\toprule
Method & Accuracy (\%) \\
\midrule
\textbf{IQC (Proposed)} & \textbf{88.07} \\
Fidelity-SWAP & 87.84 \\
Linear SVM & 90.53 \\
Logistic Regression & 90.47 \\
k-NN ($k=3$) & 92.60 \\
\bottomrule
\end{tabular}
\end{table}

To validate shot-complexity behavior, we measured classification accuracy as a function of the number of measurement shots per sample.

\begin{table}[ht]
\centering
\caption{Inference Accuracy vs. Measurement Shots: Empirical demonstration of measurement-efficient behavior for IQC relative to fidelity and variational methods.}
\begin{tabular}{l c c c c}
\toprule
Method & 10 Shots & 100 Shots & 1000 Shots & 4000 Shots \\
\midrule
IQC (Proposed) & 0.822 & 0.884 & 0.900 & 0.902 \\
Fidelity-Based & 0.654 & 0.814 & 0.906 & 0.906 \\
VQC & 0.808 & 0.808 & 0.808 & 0.808 \\
\bottomrule
\end{tabular}
\end{table}

\section{Discussion}

\subsection{Measurement Efficiency}
The IQC reaches near-peak performance with small constant shot counts, while fidelity-based and VQC models require substantially more measurements to stabilize classification scores. This reduces the duty cycle pressure on quantum hardware and improves throughput in measurement-limited settings.

The reduction in measurement shots has a direct impact on the energy consumption of the quantum-classical system. In superconducting architectures, each measurement shot involves microwave pulses and high-speed classical digitization, consuming significant power at the control electronics layer. By reducing the number of shots from $10^4$ (typical for fidelity estimation) to a small constant number of shots (often $10^1-10^2$ for IQC sign estimation), we achieve a reduction in the "energy per inference" metric. This potentially positions interference-based QML as a more energy-efficient alternative for data processing in future quantum deployments.

\subsection{Noise Robustness}
Under realistic NISQ conditions, quantum circuits are subject to decoherence and operational errors. We analyze the sensitivity of the IQC decision signal to two primary noise channels.

\subsubsection{Depolarizing Noise}
A depolarizing channel $\mathcal{E}(\rho) = (1-p)\rho + p\frac{I}{2^{n+1}}$ maps a global state $\rho$ of $n+1$ qubits to a mixture of itself and the maximally mixed state with probability $p$. For the ancilla-register system, this noise attenuates the interference signal by a factor $(1-p)^L$, where $L$ is the circuit depth:
\begin{equation}
    \langle Z_{\text{anc}} \rangle_{\text{noisy}} = (1-p)^L \mathrm{Re}\langle \phi | \psi \rangle.
\end{equation}
Crucially, since $(1-p)^L > 0$, the \textbf{sign} of the expectation value is preserved. The classification is correct as long as the shot noise does not exceed the attenuated signal.

\subsubsection{Dephasing Noise}
Dephasing or phase-damping errors $\mathcal{E}_{ph}(\rho) = (1-p)\rho + p Z\rho Z$ occur during the controlled unitary or idle times. In the Hadamard test, dephasing on the ancilla qubit adds a stochastic phase shift $\theta$ to the interference term:
\begin{equation}
    P(0) = \frac{1}{2} + \frac{1}{2} \text{Re}(e^{i\theta} \langle \phi | \psi \rangle)
\end{equation}
If the dephasing is symmetric around zero, then the expected decision score is unbiased. But if there is a systematic phase drift, then the decision boundary can be shifted. This implies that the ISDO-B is more robust to stochastic noise (such as depolarization) than to systematic calibration errors.

\subsection{Limitations of the Model}
A fundamental limitation of the current IQC framework is its reliance on a linear decision boundary in Hilbert space. While the Hilbert space dimension $2^n$ is large, separation depends heavily on the quality of the feature map $\Phi(\mathbf{x})$. If the data distribution is not linearly separable under the chosen map, the IQC behaves like a high-bias linear classifier. The dependence of inference stability on the classification margin mirrors that of classical linear classifiers and reflects intrinsic ambiguity near the decision boundary rather than a limitation of the interference-based inference mechanism.

All experiments reported in this work are conducted entirely via classical statevector simulation. No real quantum hardware experiments have been performed. Statevector simulation is noiseless by default; noise is injected via a synthetic depolarizing model. This means the reported results reflect idealized noise assumptions rather than the full complexity of real device errors (e.g., crosstalk, leakage, coherent errors). Furthermore, the current results come from a single experimental trial without multi-seed statistical validation. Real hardware validation on a backend such as IBM Quantum or IonQ, as well as multi-seed runs for mean $\pm$ std reporting, are identified as the most critical directions for strengthening the empirical claims of this work.

\subsection{Scalability Considerations}
The transition of the IQC from simulation to actual quantum hardware needs to take into account the issue of gate decomposition and ancilla usage. The Hadamard test needs controlled-unitary operations $CU$, which can be expensive in terms of circuit depth $L$ when decomposed into native gates. The interference circuit relies on a single ancilla qubit that must be connected to all qubits in the data register used in the controlled operations; restricted connectivity can introduce SWAP overhead. Readout error rates can be as high as 1--5\% in current NISQ devices, but the IQC's reliance on a single ancilla measurement helps isolate this effect. Readout error mitigation techniques can be applied to the ancilla qubit to improve the precision of the interference score without multi-qubit tomography.

The IQC framework exhibits favorable scaling compared to both classical linear models and variational quantum models. To represent a $d$-dimensional feature vector, the IQC requires $n = \lceil \log_2 d \rceil$ qubits, achieving logarithmic qubit scaling in representation (with linear gate complexity for state preparation). Once the state is prepared, the interference circuit depth $L$ is a constant or near-constant factor relative to the state preparation depth. Inference uses a single fixed circuit and a small constant number of shots to stabilize the sign decision. The classical aggregation step scales as $O(Nd)$, which is optimal for linear learning. The reduction in measurement shots also lowers the energy per inference in measurement-limited hardware settings.

\section{Conclusion and Future Work}
This paper proposes an interference-based quantum classifier that serves as a fixed inference primitive for medical image classification. By emphasizing the linear overlap $\mathrm{Re}\langle \phi | \psi \rangle$ instead of quadratic probabilities, we achieve measurement-efficient inference and noise robustness. Our method separates the learning process from the quantum computation by calculating the class centroids classically. Experimental evaluation on PCam shows that our framework is comparable to variational models while being more efficient in terms of measurement cycles and optimization costs.

Future work will extend this framework to \textbf{adaptive learning}. We will study simple update rules for the class-representative states (e.g., exponential moving averages or small ensembles) and evaluate whether they improve separability without sacrificing the measurement efficiency of the linear interference core.

\subsection{Extension to Multi-class Classification}
While this work focuses on binary classification as the core primitive, the IQC framework can be extended to multi-class problems ($K > 2$) using standard decomposition strategies such as \textbf{One-vs-All (OvA)} or \textbf{One-vs-One (OvO)}.

In the OvA approach, for $K$ classes, we learn $K$ static class-representative states $\{|\phi^{(1)}\rangle, |\phi^{(2)}\rangle, \dots, |\phi^{(K)}\rangle\}$. The interference operation is performed between the test state $|\psi\rangle$ and each class prototype. The classification rule is then:
\begin{equation}
    \hat{y} = \arg\max_c \mathrm{Re}\langle \phi^{(c)} | \psi \rangle.
\end{equation}
The shot complexity for multi-class inference scales as $O(K)$, as we must perform $K$ independent interference measurements. However, each measurement uses a small number of shots, so the overall efficiency remains significantly higher than multi-class VQCs or QSVMs. Future work will investigate more sophisticated ways to load multiple class prototypes into a single quantum memory for parallel interference, potentially reducing the scaling to $O(\log K)$.

\bibliographystyle{IEEEtran}
\begin{thebibliography}{99}

\bibitem{benedetti2019parameterized}
Marcello Benedetti, Erika Lloyd, Stefan Sack, and Mattia Fiorentini, ``Parameterized quantum circuits as machine learning models,'' \emph{Quantum Science and Technology}, vol.~4, no.~4, Art. no.~043001, 2019, doi: 10.1088/2058-9565/ab0536.

\bibitem{mcclean2018barren}
Jarrod R. McClean, Sergio Boixo, Vadim N. Smelyanskiy, Ryan Babbush, and Hartmut Neven, ``Barren plateaus in quantum neural network training landscapes,'' \emph{Nature Communications}, vol.~9, Art. no.~4812, 2018, doi: 10.1038/s41467-018-07090-4.

\bibitem{schuld2019feature}
Maria Schuld and Nathan Killoran, ``Quantum machine learning in feature Hilbert spaces,'' \emph{Physical Review Letters}, vol.~122, no.~4, Art. no.~040504, 2019, doi: 10.1103/PhysRevLett.122.040504.

\bibitem{schuld2017distance}
Maria Schuld, Mark Fingerhuth, and Francesco Petruccione, ``Implementing a distance-based classifier with a quantum interference circuit,'' \emph{EPL (Europhysics Letters)}, vol.~119, no.~6, Art. no.~60002, 2017, doi: 10.1209/0295-5075/119/60002.

\bibitem{abbas2021power}
Amira Abbas, David Sutter, Christa Zoufal, Aur\'elien Lucchi, Alessio Figalli, and Stefan W\"{o}rner, ``The power of quantum neural networks,'' \emph{Nature Computational Science}, vol.~1, no.~6, pp.~403--409, 2021, doi: 10.1038/s43588-021-00084-1.

\bibitem{havlicek2019supervised}
Vojt\v{e}ch Havl\'{\i}\v{c}ek, Antonio D. C\'orcoles, Kristan Temme, Aram W. Harrow, Abhinav Kandala, Jerry M. Chow, and Jay M. Gambetta, ``Supervised learning with quantum-enhanced feature spaces,'' \emph{Nature}, vol.~567, no.~7747, pp.~209--212, 2019, doi: 10.1038/s41586-019-0980-2.

\bibitem{sergioli2019binary}
Giuseppe Sergioli, Roberto Giuntini, and H\'ector Freytes, ``A new quantum approach to binary classification,'' \emph{PLoS ONE}, vol.~14, no.~5, e0216224, 2019, doi: 10.1371/journal.pone.0216224.

\bibitem{gonzalez2021classification}
Fabio A. Gonz\'alez, Vladimir Vargas-Calder\'on, and Herbert Vinck-Posada, ``Classification with quantum measurements,'' arXiv:2004.01227v2, 2021.

\bibitem{blank2020kernel}
Carsten Blank, Daniel K. Park, June-Koo Kevin Rhee, and Francesco Petruccione, ``Quantum classifier with tailored quantum kernel,'' \emph{npj Quantum Information}, vol.~6, Art. no.~41, 2020, doi: 10.1038/s41534-020-0272-6.

\bibitem{blank2022compact}
Carsten Blank, Adenilton J. da Silva, Lucas P. de Albuquerque, Francesco Petruccione, and Daniel K. Park, ``Compact quantum kernel-based binary classifier,'' \emph{Quantum Science and Technology}, vol.~7, no.~4, Art. no.~045007, 2022, doi: 10.1088/2058-9565/ac842c.

\bibitem{schuld2018book}
Maria Schuld and Francesco Petruccione, \emph{Supervised Learning with Quantum Computers}. Springer, 2018.

\bibitem{plesch2011quantum}
Martin Plesch and \v{C}aslav Brukner, ``Quantum-state preparation with universal gate sets,'' \emph{Physical Review A}, vol.~83, no.~3, Art. no.~032302, 2011, doi: 10.1103/PhysRevA.83.032302.

\bibitem{pcam_01}
Reek Majumdar, Biswaraj Baral, Bhavika Bhalgamiya, and Taposh Dutta Roy, ``Histopathological cancer detection using hybrid quantum computing,'' arXiv:2302.04633v1, 2023.

\bibitem{pcam_03}
Nandika Goyal, ``Quantum enhanced histopathologic cancer detection,'' M.S. thesis, Arizona State University, 2024.

\bibitem{pcam_04}
Raghavendra Selvan and Erik B. Dam, ``Tensor networks for medical image classification,'' in \emph{Proc. Medical Imaging with Deep Learning (MIDL)}, PMLR, vol.~121, pp.~721--731, 2020.

\bibitem{pcam_veeling2018}
Bastiaan S. Veeling, Jasper Linmans, Jim Winkens, Taco Cohen, and Max Welling, ``Rotation equivariant CNNs for digital pathology,'' in \emph{Medical Image Computing and Computer Assisted Intervention -- MICCAI 2018}, 21st International Conference, Granada, Spain, Proceedings, Part II, pp.~210--218. Springer International Publishing, 2018.

\bibitem{cover1965geometrical}
Thomas M. Cover, ``Geometrical and statistical properties of systems of linear inequalities with applications in pattern recognition,'' \emph{IEEE Transactions on Electronic Computers}, vol.~EC-14, no.~3, pp.~326--334, 1965, doi: 10.1109/PGEC.1965.264137.

\bibitem{ekert2002direct}
Artur K. Ekert, Carlos M. Alves, Daniel K. L. Oi, M\"{a}rcin Horodecki, Pawe\"{l} Horodecki, and L. C. Kwek, ``Direct estimations of linear and nonlinear functionals of a quantum state,'' \emph{Physical Review Letters}, vol.~88, no.~21, Art. no.~217901, 2002, doi: 10.1103/PhysRevLett.88.217901.

\end{thebibliography}

\end{document}
